%! Introduction to Functions and Change — Unit 1, Lesson 1.0 — Slides (Beamer)
\documentclass[aspectratio=169, 11pt]{beamer}
\usepackage{precalculus-beamer}   % provides \plumheader{} and \sectionlabel[color]{}

\begin{document}

% TITLE SLIDE
{
\setbeamercolor{background canvas}{bg=plum}
\begin{frame}[plain]
  \vspace{2.8cm}\hspace{0.55cm}%
  \begin{minipage}{0.9\textwidth}
    {\color{goldacc}\bfseries\small LESSON 1.0}\\[0.5cm]
    {\color{white}\bfseries\LARGE Introduction to Functions and Change}\\[1.2cm]
    {\color{lilacmid}\small Precalculus~$\cdot$~Unit 1: Functions and Change}
  \end{minipage}
\end{frame}
}

% HOOK SLIDE
\begin{frame}[t, plain]
  \plumheader{Hook: The Trustworthy Machine}
  \sectionlabel[goldacc]{Think About It}

  A vending machine and a claw machine both take your money.

  \bigskip
  \begin{columns}[T]
    \column{0.48\textwidth}
      \begin{block}{Vending Machine}
        \begin{itemize}\setlength{\itemsep}{6pt}
          \item Press \textbf{B4} \(\to\) same bag of chips
          \item Every. Single. Time.
          \item Same input, same output
        \end{itemize}
      \end{block}

    \column{0.48\textwidth}
      \begin{block}{Claw Machine}
        \begin{itemize}\setlength{\itemsep}{6pt}
          \item Drop the claw \(\to\) who knows?
          \item Same move, different results
          \item Same input, \textbf{unpredictable} output
        \end{itemize}
      \end{block}
  \end{columns}

  \bigskip
  \textit{Discussion:} Which machine can you \textbf{trust}? What would it mean for
  the vending machine to ``break the rules''?
\end{frame}

% ROW 1 — INPUT & OUTPUT
\begin{frame}[t, plain]
  \plumheader{Input \& Output}
  \sectionlabel[goldacc]{Notes row INPUT \& OUTPUT $\cdot$ I Do: problem 1 $\cdot$ We Do: problems 2--4}

  \begin{block}{Two quantities, one responds}
    The quantity we \textbf{choose or track} is the \textbf{input} (independent
    variable). The quantity that \textbf{responds} is the \textbf{output} (dependent
    variable). Ordered pairs are (input, output).
  \end{block}

  \medskip
  \begin{center}
  \begin{tikzpicture}[>=stealth]
    \node[draw=plum, very thick, fill=lilac, rounded corners=3pt, minimum width=3.4cm,
          minimum height=1.1cm, align=center, font=\small] (box) at (4.3,0)
          {\textbf{pay-per-minute phone plan}};
    \node[font=\small] (in) at (0,0) {minutes talked};
    \node[font=\small] (out) at (8.8,0) {cost of the bill};
    \draw[->, very thick, plum] (in) -- (box);
    \draw[->, very thick, plum] (box) -- (out);
    \node[font=\scriptsize\itshape, below] at (in.south) {you choose: input};
    \node[font=\scriptsize\itshape, below] at (out.south) {it responds: output};
  \end{tikzpicture}
  \end{center}

  \medskip
  \textbf{The ``you-control'' test:} which quantity do \textit{you} choose? That one is
  the input.

  \smallskip
  \begin{enumerate}\setlength{\itemsep}{2pt}
    \item A lit candle burning down
    \item A ball dropped from a rooftop
    \item Pumping gas into a car --- plus one ordered pair
    \item ``The phone bill decides how many minutes you talk.'' What is backward?
  \end{enumerate}
\end{frame}

% ROW 2 — FUNCTION
\begin{frame}[t, plain]
  \plumheader{The Dependable Machine: Function}
  \sectionlabel[goldacc]{Notes row FUNCTION $\cdot$ I Do: problem 5 $\cdot$ We Do: problems 6--8}

  \begin{block}{Definition (informal)}
    A rule is a \textbf{function} when each input produces \textbf{exactly one}
    output. In an arrow diagram: \textbf{exactly one arrow leaves each input}.
  \end{block}

  \begin{columns}[T]
    \column{0.48\textwidth}
      \begin{center}
      \begin{tikzpicture}[>=stealth, font=\small]
        \foreach \n/\y in {A1/1.1, B4/0.4, C2/-0.3} \node (i\n) at (0,\y) {\n};
        \foreach \n/\y/\k in {chips/1.1/a, pretzels/0.4/b, gum/-0.3/c} \node (o\k) at (2.6,\y) {\n};
        \draw[->, thick, plum] (iA1) -- (oa);
        \draw[->, thick, plum] (iB4) -- (ob);
        \draw[->, thick, plum] (iC2) -- (oc);
        \node[font=\small\bfseries\color{plum}] at (1.3,-1.0) {vending: a function};
      \end{tikzpicture}
      \end{center}
    \column{0.48\textwidth}
      \begin{center}
      \begin{tikzpicture}[>=stealth, font=\small]
        \node (d1) at (0,0.8) {drop 1};
        \node (d2) at (0,-0.1) {drop 2};
        \foreach \n/\y/\k in {bear/1.1/a, nothing/0.4/b, ball/-0.3/c} \node (p\k) at (2.6,\y) {\n};
        \draw[->, thick, redacc] (d1) -- (pa);
        \draw[->, thick, redacc] (d1) -- (pb);
        \draw[->, thick, redacc] (d2) -- (pc);
        \node[font=\small\bfseries\color{redacc}] at (1.3,-1.0) {claw: not a function};
      \end{tikzpicture}
      \end{center}
  \end{columns}

  \smallskip
  \begin{columns}[T]
    \column{0.48\textwidth}
      \begin{enumerate}\setlength{\itemsep}{1pt}
        \setcounter{enumi}{4}
        \item Student \(\to\) their ID number
        \item Month \(\to\) students born then
      \end{enumerate}
    \column{0.48\textwidth}
      \begin{enumerate}\setlength{\itemsep}{1pt}
        \setcounter{enumi}{6}
        \item Number \(\to\) its double
        \item Shoe size \(\to\) a student who wears it
      \end{enumerate}
  \end{columns}
\end{frame}

% ROW 3 — MANY-TO-ONE (THE CRUX)
\begin{frame}[t, plain]
  \plumheader{The Fine Print: Many-to-One vs.\ One-to-Many}
  \sectionlabel[redacc]{Notes row MANY-TO-ONE $\cdot$ I Do: 9 $\cdot$ We Do: 10--12 (crux: 11, 12)}

  \begin{columns}[T]
    \column{0.48\textwidth}
      \begin{block}{Allowed: many-to-one}
        Two inputs \textbf{share} an output. Twins Ana and Ben \(\to\) May 3.
        Still a function.
      \end{block}
    \column{0.48\textwidth}
      \begin{block}{Broken: one-to-many}
        One input, \textbf{two outputs}. Shoe size 9 \(\to\) Dev and Eli.
        Not a function.
      \end{block}
  \end{columns}

  \medskip
  \begin{center}
    \Large\textbf{Count the arrows \textcolor{plum}{leaving each input} ---\\
    never the arrows arriving at an output.}
  \end{center}

  \begin{enumerate}\setlength{\itemsep}{2pt}
    \setcounter{enumi}{8}
    \item \((1, 5),\ (2, 5),\ (3, 7)\) \qquad 10.\ \((4, 1),\ (4, 2),\ (6, 3)\)
    \setcounter{enumi}{10}
    \item Two students both scored 80. Maya: ``So student \(\to\) score is not a function.''
    \item Student \(\to\) birthday is a function. Is birthday \(\to\) student?
  \end{enumerate}
\end{frame}

% ROW 4 — READING GRAPHS
\begin{frame}[t, plain]
  \plumheader{A Graph Tells a Story}
  \sectionlabel[goldacc]{Notes row READING GRAPHS $\cdot$ I Do: problem 13 $\cdot$ We Do: problems 14--15}

  \begin{columns}[T]
    \column{0.5\textwidth}
      \begin{center}
      \begin{tikzpicture}[x=0.5cm, y=0.42cm]
        \draw[linegray, very thin, step=1] (0,0) grid (10,6);
        \draw[->, thick] (0,0) -- (10.6,0) node[below right] {\small time (min)};
        \draw[->, thick] (0,0) -- (0,6.6) node[above] {\small distance (blocks)};
        \foreach \x in {2,4,6,8,10} \node[below] at (\x,0) {\small \x};
        \foreach \y in {1,3,5} \node[left] at (0,\y) {\small \y};
        \draw[very thick, plum] (0,0) -- (4,5) node[midway, above left] {\small\bfseries A};
        \draw[very thick, plum] (4,5) -- (6,5) node[midway, above] {\small\bfseries B};
        \draw[very thick, plum] (6,5) -- (10,0) node[midway, above right] {\small\bfseries C};
      \end{tikzpicture}
      \end{center}

    \column{0.46\textwidth}
      Read it left to right, like a story:
      \begin{itemize}\setlength{\itemsep}{5pt}
        \item \textbf{A}: rising --- \textbf{increasing} --- riding away
        \item \textbf{B}: flat --- \textbf{constant} --- resting
        \item \textbf{C}: falling --- \textbf{decreasing} --- riding home
      \end{itemize}
      A flat stretch is not ``nothing happening'': time keeps moving; the
      \textbf{output} holds steady.
  \end{columns}

  \medskip
  \begin{enumerate}\setlength{\itemsep}{2pt}
    \setcounter{enumi}{12}
    \item How far from home at \(t = 5\)? What is happening on B?
    \item When is the distance decreasing? What is the rider doing?
    \item Is distance from home a function of time? Explain.
  \end{enumerate}
\end{frame}

% YOU DO
\begin{frame}[t, plain]
  \plumheader{You Do --- Tables \& Flips}
  \sectionlabel[redacc]{Notes row TABLES \& FLIPS $\cdot$ problems 16--19, alone, 14 minutes}

  \begin{columns}[T]
    \column{0.48\textwidth}
      \begin{tabular}{|c|c|c|c|c|c|}
      \hline Hours & 0 & 1 & 2 & 3 & 4 \\ \hline Height (cm) & 20 & 17 & 14 & 14 & 11 \\ \hline
      \end{tabular}
    \column{0.48\textwidth}
      \begin{tabular}{|c|c|c|c|c|}
      \hline Hours studied & 0 & 1 & 2 & 3 \\ \hline Quiz score & 60 & 70 & 80 & 80 \\ \hline
      \end{tabular}
  \end{columns}

  \medskip
  \begin{enumerate}\setlength{\itemsep}{6pt}
    \setcounter{enumi}{15}
    \item The candle: when is the height decreasing, and when constant? Which bike-graph
          segment matches?
    \item The quiz: inputs 2 and 3 share 80. Still a function of hours studied?
    \item A runner jogs out to a bench, rests, and jogs back --- at 0.3\,km once going out
          and once coming back. Is distance a function of time?
    \item Flip it: is \textit{time} a function of \textit{distance}?
  \end{enumerate}
\end{frame}

% DEBRIEF
\begin{frame}[t, plain]
  \plumheader{Debrief --- The Headline}
  \sectionlabel[goldacc]{Share out 16--19, one answer each $\cdot$ then say it back}

  \bigskip
  \begin{block}{A function}
    sends each \textbf{input} to \textbf{exactly one} output.\\[4pt]
    Two inputs \textbf{may} share an output. One input may \textbf{never} have two.
  \end{block}

  \bigskip
  \textbf{Where this unit goes:} 1.1 gives functions their notation --- \(f(x)\).
  Everything after it asks how one quantity changes with respect to another. See the
  lessons table on your unit cover.
\end{frame}

% HOMEWORK LAUNCH
\begin{frame}[t, plain]
  \plumheader{Start Your Homework}
  \sectionlabel[redacc]{Homework Launch $\cdot$ 3 minutes, alone, silent}

  \begin{itemize}\setlength{\itemsep}{8pt}
    \item Copy the due date into the \textbf{Homework} row of your cover sheet:
          \quad Due: \rule{3cm}{0.4pt}
    \item Start \textbf{problem 1} and \textbf{problem 7} now. For every function
          question, count the arrows leaving each \textit{input}.
    \item The \textbf{Homework} (last two pages) is graded.
    \item \textbf{AP Practice} (the two pages before it) is \textbf{extra credit}
          --- optional, AP-exam style. Do the homework first.
  \end{itemize}
\end{frame}

\end{document}
